\section{Human AI Interaction}

\begin{ejercicio}
    Explain the three aspects (ethically-aligned design, human factors design, and technology enhancement) of Human-AI Interaction. Identify which aspect is neglected in the given scenario and justify your answer.
    \begin{enumerate}
        \item A self-driving car is designed to navigate through traffic and make decisions based on real-time data. However, it lacks a user-friendly interface for passengers to understand the car's decision-making process.
    \end{enumerate}

    The three aspects of Human-AI Interaction are:
    \begin{enumerate}
        \item \ul{Ethically-aligned design}: This aspect focuses on ensuring that AI systems are designed with ethical considerations in mind, such as fairness, transparency, and accountability. It involves creating AI systems that align with human values and societal norms.
        \item \ul{Human factors design}: This aspect emphasizes the importance of designing AI systems that are user-friendly and take into account human cognitive and physical capabilities. It involves creating interfaces and interactions that are intuitive and easy for humans to understand and use.
        \item \ul{Technology enhancement}: This aspect involves leveraging AI technology to enhance human capabilities and improve decision-making processes. It focuses on how AI can augment human abilities and provide valuable insights to support human decision-making.
    \end{enumerate}

    The neglected aspect in the given scenario is:
    \begin{enumerate}
        \item In the scenario of the self-driving car, the neglected aspect is \ul{human factors design}. The car lacks a user-friendly interface for passengers to understand the car's decision-making process. This means that while the car may be technologically advanced and ethically aligned, it fails to consider the human experience and usability.
    \end{enumerate}
\end{ejercicio}

\begin{ejercicio}
    Explain the concept of ``human-in-the-loop'' and its importance in AI systems. Provide an example of a situation where human-in-the-loop is crucial.

    The concept of ``human-in-the-loop'' refers to the integration of human oversight and intervention in AI systems. The human only provides feedback and guidance to the AI system, leading to an improved decision-making process. This approach is important because it allows for human judgment, ethical considerations, and contextual understanding to be incorporated into AI decision-making, which can help prevent errors and biases that may arise from fully automated systems.

    An example of a situation where human-in-the-loop is crucial is in medical diagnosis. AI systems can analyze medical images and provide diagnostic suggestions, but the final decision should be made by a human doctor who can consider the patient's medical history, symptoms, and other relevant factors. This ensures that the diagnosis is accurate and takes into account the nuances of individual cases.
\end{ejercicio}

\begin{ejercicio}
    Is there always a trade-off between human control and the automatization of a system? Explain how this two aspects are related.

    There is not always a trade-off between human control and the automatization of a system. While it is true that increasing automation can reduce the level of human control, it is possible to design systems that balance both aspects. Therefore, a two dimensional approach can be taken, where totally automated systems are at the right end of the spectrum and fully human-controlled systems are at the upper end of the spectrum.
    \begin{itemize}
        \item Some aspects would have no human control and be fully automated, such as the braking system of a car, which is designed to respond automatically in emergency situations.
        \item Other aspects would have a high level of human control and be less automated, such as playing music.
        \item Some aspects provide both human control and automation, such as an elevator or a calculator, which can operate automatically but also allow for human input and control. These would be in the upper right quadrant of the spectrum.
    \end{itemize}
\end{ejercicio}

\begin{ejercicio}
    Name two forms of explanations and describe them. Why are explanations important? Give the four reasons that make explanations important.

    Some forms of explanations include:
    \begin{enumerate}
        \item \ul{Visual explanations}: These explanations use visual representations, such as hot maps or saliency maps, to illustrate how an AI model makes decisions. They help users understand which features or inputs are most influential in the model's predictions. For instance, when detecting if a patient has a tumor, a visual explanation can highlight the specific areas of the medical image that contributed to the AI's diagnosis.
        \item \ul{Rule-based explanations}: These explanations provide a set of rules or conditions that describe how the AI model arrives at its decisions. They can be presented in a human-readable format, allowing users to follow the logic behind the model's predictions. For example, a rule-based explanation for a loan approval AI system might state that if an applicant has a credit score above 700 and a stable income, they are likely to be approved.
    \end{enumerate}

    Explanations are important for several reasons:
    \begin{enumerate}
        \item \ul{Trust}: Explanations help build trust between users and AI systems by providing transparency into how decisions are made.
        \item \ul{Accountability}: Explanations allow users to hold AI systems accountable for their decisions, as they can understand the reasoning behind those decisions.
        \item \ul{Debugging and improvement}: Explanations can help developers identify errors or biases in AI models, leading to improvements and refinements in the system.
        \item \ul{Ethical considerations}: Explanations ensure that AI systems operate in an ethically responsible manner, as users can assess whether the decisions align with societal values and norms.
    \end{enumerate}
\end{ejercicio}

\begin{ejercicio}
    We came to the conclusion that AI should be explainable. What are the 4 steps to explainability?
    \begin{enumerate}
        \item \ul{Identify the audience}: Determine who will be receiving the explanation and tailor the explanation to their level of understanding and expertise. For instance, a technical explanation may be suitable for data scientists, while a simplified explanation may be more appropriate for end-users.
        \item \ul{Determine the purpose of the explanation}: Clarify why the explanation is being provided and what information needs to be conveyed. This could include providing insights into the decision-making process, addressing user concerns, or ensuring compliance with regulations.
        \item \ul{Choose the content of the explanation}: Decide what information will be included in the explanation, such as the model's inputs, outputs, and reasoning process. This may involve selecting relevant features, highlighting important factors, or providing examples to illustrate the model's behavior.
        \item \ul{Select the explanation modality}: Choose the appropriate format or method for delivering the explanation, such as visualizations, textual descriptions, or interactive interfaces. The chosen modality should effectively communicate the information to the intended audience and facilitate understanding.
    \end{enumerate}
\end{ejercicio}

\begin{ejercicio}
    A doctor receives a diagnosis suggestion from an AI but doesn't follow it. How and why can explanation modalities be used to increase trust and understanding?\\

    In this case, the phenomenon that happens is called ``underreliance''. Explanation modalities can be used to increase trust and understanding by providing the doctor with clear and transparent information about how the AI arrived at its diagnosis suggestion. By using visual explanations, such as highlighting relevant features in medical images, or rule-based explanations that outline the decision-making process, the doctor can gain insights into the AI's reasoning. This can help the doctor assess the validity of the AI's suggestion, identify any potential errors or biases, and ultimately make a more informed decision. By enhancing transparency and providing a deeper understanding of the AI's capabilities, explanation modalities can foster trust and confidence in the AI system.
\end{ejercicio}

\begin{ejercicio}
    Briefly describe the thought experiment ``Chinese room''. What can you deduce about the intelligence of programs from this experiment?

    The ``Chinese room'' thought experiment is designed to challenge the notion that a computer program can possess true understanding or intelligence. In this experiment, a person who does not understand Chinese is placed in a room with a set of rules and symbols that allow them to manipulate Chinese characters. The person receives Chinese input and, according to the rules provided, produces appropriate Chinese output without understanding the meaning of the characters.

    From this experiment, we can deduce that a program may be able to simulate understanding or intelligence without actually possessing it. The program can produce appropriate responses based on rules and symbols, but it does not have genuine comprehension or consciousness.
\end{ejercicio}